\scanpage{90}
[ Si \inserted{petites} $\underrightarrow{\lim}$ existent \inserted{ds $\unsure{E|\mathcal E}_b$,} et les $f^*$ y commutent :

\[
L\xrightarrow{i}\underline{\Gamma}(E/B)\qquad\text{a un adjoint à gauche }\inserted{i_*}
\]

Si petites $\underleftarrow{\lim}$ existent ds $E_b$, et les $f^*$ y commutent,

\[
L^\circ\xrightarrow{j^\circ}\underline{\Gamma}(E^{\triangle}/B)
\]

a un adjoint à gauche, donc

\[
L\mathrel{\overset{j}{\unsure{\longrightarrow|\rule{1.5em}{.35pt}}}}\underline{\Gamma}(E^\triangle/B)^\circ
\]

a un adjoint à droite $j_+$

Si on cherche des générateurs de $L$, il y a
lieu d'utiliser $i_*$, car \deleted{pour les} si on a des
gén $X_\lambda$ de $\underline{\Gamma}(E/B)$, comme

\[
\mathrm{Hom}_{\underline{\Gamma}}(X_\lambda,i(Y))\simeq
\mathrm{Hom}_L(i_*(X_\lambda),Y),
\]

on voit que la famille des $i_*(X_\lambda)$ forme un
syst. de gén. de $L$. Or, pour avoir un \inserted{(\uline{petit})}
syst. de gén de $\underline{\Gamma}$, le cor. 2 \inserted{de Deligne} appliqué
à $E^\circ\mathbin{\unsure{-|\to}}B^\circ$ nous dit qu'il suffit que
$p$ co\uline{fibrant} et que les petites sommes
existent ds les $E_b$, et que dans ceux-ci
il y ait des \deleted{gén.} petits ens. gén.

\[
\left|\ \begin{aligned}
&\text{Petites }\underrightarrow{\mathrm{Lim}}\ \text{existent ds les }E_b,\ \text{les }f^*\text{ y commutent}\\
&\text{Petit ens de générateurs dans les }E_b
\end{aligned}\right.
\]

% archive sheet number: 89
